\begin{enumerate}
\item[\textup{\SGAInlineRefTarget{sga3:xix:proposition:4:12:item:c}(c)}]
\textit{Suppose that}
\[
  (\alpha'\circ f)_{\bar s_0}=q\alpha_{\bar s_0},
\]
\textit{where \(q>0\) is an integer such that
\(x\longmapsto x^q\) is an endomorphism of
\(\mathbb G_{a,S}\).  There then exist an open subset \(U\subset S\)
containing \(s_0\) and a unique morphism of
\(\mathcal O_S\)-modules}
\[
  h\colon L^{\otimes q}|_U\longrightarrow L'|_U
\]
\textit{such that \(g_U\) is the composite}
\[
  \mathbb W(L)_U
  \xrightarrow{\ x\mapsto x^q\ }
  \mathbb W(L^{\otimes q})_U
  \xrightarrow{\ \mathbb W(h)\ }
  \mathbb W(L')_U.
\]
\end{enumerate}

\SGAProof{Proof of \SGARefLink{sga3:xix:proposition:4:12}{Proposition 4.12}.}
Let us prove \textup{(a)}.  Since the conclusion is local on \(S\),
we may suppose that
\(\mathbb W(L)=\mathbb W(L')=\mathbb G_{a,S}\), and hence that \(g\)
is given by a polynomial
\[
  g(X)=\sum_{n\geq0}a_nX^n,
  \qquad
  a_n\in\Gamma(S,\mathcal O_S).
\]
Condition \textup{\((\star)\)}, which relates \(f\) and \(g\), is the
identity in \(\Gamma(S',\mathcal O_{S'})[X]\)
\[
  \sum_{n\geq0}a_n\alpha'\bigl(f(t)\bigr)X^n
  =
  \sum_{n\geq0}a_n\alpha(t)^nX^n.
\]
Thus, for every \(n\geq0\), every \(S'\to S\), and every
\(t\in T(S')\),
\[
  a_n\bigl(\alpha'(f(t))-\alpha(t)^n\bigr)=0.
\]

For every \(n\geq0\), let \(S_n\) be the set of points \(s\in S\)
such that
\[
  (\alpha'\circ f)_{\bar s}=n\alpha_{\bar s}.
\]
It is known \textup{(\SGARefLink{sga3:ix:corollary:5:3}{Exposé IX, 5.3})} that the \(S_n\) are both open
and closed and that
\[
  (\alpha'\circ f)_{S_n}=n\alpha_{S_n}.
\]
Moreover, since \(\alpha_{\bar s_0}\ne0\), after shrinking \(S\) we
may suppose that \(\alpha\) is nonzero on every fiber \textup{(same
reference)}; it follows that the \(S_n\) are pairwise disjoint.
After shrinking \(S\) once more, we may therefore suppose that one
of the following two cases holds: \(S=S_n\) for some \(n\), or every
\(S_n\) is empty.

Let \(m\geq0\) be such that \(S_m=\varnothing\).  We claim that
\(a_m=0\).  Indeed, \(\alpha'\circ f\) and \(m\alpha\) are distinct
on every fiber of \(S\), and we have:

\par\medskip\noindent
\textbf{\textit{\SGARefLink{sga3:xix:lemma:4:13}{Lemma 4.13}.}}---
\textit{Let \(S\) be a scheme, let \(T\) be an \(S\)-torus, and let
\(\alpha,\alpha'\) be two characters of \(T\) which are distinct on
every fiber.  There exist an fpqc covering family
\(\{S_i\to S\}\) and, for every \(i\), an element \(t_i\in T(S_i)\)
such that}
\[
  \alpha(t_i)-\alpha'(t_i)=1.
\]
\SGARefTarget{sga3:xix:lemma:4:13}

The lemma is immediate by reduction first to the diagonalizable
case, and then to \(T=\mathbb G_{m,S}\).

We return to the proof of \SGARefLink{sga3:xix:proposition:4:12}{Proposition 4.12}.  We have just proved
\textup{(a)}.  In cases \textup{(b)} and \textup{\SGARefLink{sga3:xix:proposition:4:12:item:c}{(c)}}, there is an
integer \(n\) such that \(S=S_n\) \textup{(\(n=q\) in
\textup{\SGARefLink{sga3:xix:proposition:4:12:item:c}{(c)}})}.  The preceding argument gives \(a_m=0\) for
\(m\ne n\), and therefore
\[
  g(X)=a_nX^n,
  \qquad
  a_n\in\Gamma(S,\mathcal O_S).
\]
This proves \textup{(c)} at once.  In case \textup{(b)}, we know
that \(a_n(s_0)\ne0\); we may therefore suppose that \(a_n\) is
invertible on \(S\).  It follows that \(x\longmapsto x^n\) is an
endomorphism of \(\mathbb G_{a,S}\), by the hypothesis in
\textup{(b)}, and this completes the proof.

